% !PRINTED 57
\weberpage{57}
%<<CONT

not without leaving different powers of $t$ in the separate terms.

By increasing the number of variables one can transform every non-homogeneous function into a homogeneous one of the same degree. If $m-1$ is the number of variables in a non-homogeneous function of degree $n$, then we put
\[
 x={x_1\over x_m},\qquad y={x_2\over x_m},\qquad z={x_3\over x_m},\ldots,
\]
and obtain in
\[
 x_m^n F\!\left({x_1\over x_m},{x_2\over x_m},{x_3\over x_m},\ldots\right)
\]
a whole homogeneous function of degree $n$ in the variables $x_1,x_2,\ldots,x_m$, which we denote by
\[
 \Phi(x_1,x_2,\ldots,x_m).
\]

It is sometimes advisable to write homogeneous functions of several variables \gesperrt{with the polynomial coefficients}. We therefore put
\begin{equation}
 \Phi(x_1,x_2,\ldots,x_m)=
 \sum{\Pi(n)\over\Pi(\alpha_1)\Pi(\alpha_2)\cdots\Pi(\alpha_m)}
 A_{\alpha_1,\alpha_2,\ldots,\alpha_m}
 x_1^{\alpha_1}x_2^{\alpha_2}\cdots x_m^{\alpha_m},
\tag{2}
\end{equation}
where the sum is extended over all non-negative numbers satisfying the condition
\begin{equation}
 \alpha_1+\alpha_2+\cdots+\alpha_m=n.
\tag{3}
\end{equation}
This notation, without the restriction (3), is also applicable to non-homogeneous functions.

But the homogeneous function can also be represented as
\begin{equation}
 \Phi(x_1,x_2,\ldots,x_m)=
 \sum A_{\nu_1,\nu_2,\ldots,\nu_n}x_{\nu_1}x_{\nu_2}\cdots x_{\nu_n},
\tag{4}
\end{equation}
where each of the indices $\nu_1,\nu_2,\ldots,\nu_n$, independently of the others, runs through the range of values $1,2,\ldots,m$. Thus the sum (4) consists of $m^n$ terms, but not all of them are different from one another. The product $x_{\nu_1}x_{\nu_2}\cdots x_{\nu_n}$ remains unchanged if the indices $\nu_1,\nu_2,\ldots,\nu_n$ are permuted among themselves in any way. The number of permutations of $n$ elements is $\Pi(n)$. If among these elements $\alpha_1,\alpha_2,\ldots$ respectively are equal to one another, then the number of permutations is reduced to
\[
 {\Pi(n)\over \Pi(\alpha_1)\Pi(\alpha_2)\cdots},
\]
%CONT>>
